\section*{Chapter \ref{ch:b2:linalg} --- Linear Algebra}

\begin{solution}{exo:b2:linalg:1}
Three forms in a $3$-dimensional dual: freeness suffices. A relation
$\alpha\varphi_1 + \beta\varphi_2 + \gamma\varphi_3 = 0$ evaluated at
$(1,0,0), (0,1,0), (0,0,1)$ gives $\alpha + \gamma = 0$, $\alpha +
\beta = 0$, $\beta + \gamma = 0$, whence $\alpha = \beta = \gamma =
0$.

Pre-dual basis $(u_1, u_2, u_3)$: solve $\varphi_i(u_j) =
\delta_{ij}$. Writing $u_j = (x, y, z)$: for $u_1$: $x + y = 1$,
$y + z = 0$, $x + z = 0$ gives $u_1 = \bigl(\tfrac12, \tfrac12,
-\tfrac12\bigr)$; symmetrically $u_2 = \bigl(-\tfrac12, \tfrac12,
\tfrac12\bigr)$, $u_3 = \bigl(\tfrac12, -\tfrac12, \tfrac12\bigr)$.
\end{solution}

\begin{solution}{exo:b2:linalg:2}
First matrix: the only nonzero-product permutation sends $1 \mapsto
3$, $2 \mapsto 2$, $3 \mapsto 1$ --- the transposition $(1\,3)$,
signature $-1$: determinant $-abc$.

Second: a permutation with nonzero product cannot mix the two
blocks (an entry linking them is $0$), so it splits as a permutation
of $\{1,2\}$ times one of $\{3,4\}$, and the signature is the
product of the two signatures: the sum factorizes as
\[
(ad - bc)(eh - fg) .
\]
General rule suggested (and true, same proof): the determinant of a
block-diagonal matrix is the product of the determinants of the
blocks.
\end{solution}

\begin{solution}{exo:b2:linalg:3}
$F$ is defined by the two independent equations $\varphi_1(x,y,z,t)
= x + y - z - t = 0$ and $\varphi_2 = x - 2y = 0$: by
\cref{thm:b2:linalg:annihilator} read backwards, $F^\circ =
\operatorname{Vect}(\varphi_1, \varphi_2)$ --- they lie in $F^\circ$
by construction, they are free (not proportional), and $\dim F^\circ
= 4 - \dim F = 4 - 2 = 2$ since $\dim F = 2$ (two independent
equations in $\R^4$). Basis: $(\varphi_1, \varphi_2)$; dimensions:
$2 + 2 = 4$, as the theorem demands.
\end{solution}

\begin{solution}{exo:b2:linalg:4}
The $\varphi_i$ are $n + 1$ forms on an $(n+1)$-dimensional space:
freeness suffices. If $\sum_i \lambda_i \varphi_i = 0$, evaluate on
the Lagrange polynomial $L_j$ of the nodes: $\lambda_j = 0$. The
pre-dual basis is $(L_0, \dots, L_n)$, since $\varphi_i(L_j) =
L_j(a_i) = \delta_{ij}$.

For the integral form with nodes $0, \frac12, 1$ on $\R_2[X]$:
$\int_0^1 P = \sum_i c_i P(a_i)$ with $c_i = \int_0^1 L_i$. Compute:
$L_0 = 2(X - \tfrac12)(X - 1)$, $\int_0^1 L_0 = \frac16$; $L_1 =
-4X(X-1)$, $\int_0^1 L_1 = \frac46$; $L_2 = 2X(X - \tfrac12)$,
$\int_0^1 L_2 = \frac16$. Hence
\[
\int_0^1 P = \frac{1}{6}\Bigl(P(0) + 4P\bigl(\tfrac12\bigr) +
P(1)\Bigr)
\quad (P \in \R_2[X]) :
\]
Simpson's rule, exact on quadratics --- a statement about dual
bases.
\end{solution}

\begin{solution}{exo:b2:linalg:5}
$\operatorname{im} u = Ka$ for some $a \neq 0$; then $u(x) =
\varphi(x)\,a$ where $\varphi(x)$ is the coordinate of $u(x)$ on $a$
--- linear in $x$. Trace: complete $a = e_1$ into a basis; the
matrix of $u$ has columns $\varphi(e_j)\,e_1$, so its only diagonal
entry is $\varphi(e_1) = \varphi(a)$: $\operatorname{tr} u =
\varphi(a)$. Then
\[
u^2(x) = \varphi(x)\, u(a) = \varphi(x)\varphi(a)\, a
= (\operatorname{tr} u)\, u(x).
\]
Determinant, in two cases. \emph{If $\varphi(a) \neq 0$:} take any
basis of the hyperplane $\ker\varphi$ and append $a$. Then $u$ kills
$\ker\varphi$ (there $u(x) = \varphi(x)a = 0$) and $u(a) =
\varphi(a)\,a$: the matrix of $I + u$ is diagonal, $(1, \dots, 1,\,
1 + \varphi(a))$, so $\det(I + u) = 1 + \varphi(a) = 1 +
\operatorname{tr} u$. \emph{If $\varphi(a) = 0$:} then $a \in
\ker\varphi$; take a basis of $\ker\varphi$ whose first vector is
$a$, and append a vector $b$ with $\varphi(b) = 1$. Then $I + u$
fixes the basis of $\ker\varphi$ and sends $b \mapsto b + a$:
triangular with unit diagonal, $\det(I + u) = 1 = 1 +
\operatorname{tr} u$. Both cases agree with the formula.
\end{solution}

\begin{solution}{exo:b2:linalg:6}
By \cref{exo:b2:linalg:9}, the hyperplane is $H_A = \{M :
\operatorname{tr}(AM) = 0\}$ with $A \neq 0$.

\emph{If $A = \lambda I$:} $H_A$ is the zero-trace hyperplane; the
matrix of the $n$-cycle permutation (ones in positions $(i, i+1)$
and $(n, 1)$) is invertible (its determinant is $\pm 1$ by
\cref{ex:b2:linalg:permexample}'s computation) and has zero trace.

\emph{If $A$ is not scalar:} first find an invertible $P$ such that
$B = P^{-1}AP$ has a nonzero off-diagonal entry $b_{ji}$ ($j \neq
i$). Indeed, if $A$ already has one, take $P = I$; if $A$ is
diagonal with two distinct entries $d_1 \neq d_2$, conjugating by
the transvection $P = I + E_{12}$ produces the off-diagonal entry
$d_1 - d_2 \neq 0$ (compute: $P^{-1}AP = A + (d_1 - d_2)E_{12}$);
and a diagonal matrix with all entries equal is
scalar, excluded. Now set $M' = I + tE_{ij}$ with $t =
-\operatorname{tr}(B)/b_{ji}$: then
\[
\operatorname{tr}(BM') = \operatorname{tr} B + t\,b_{ji} = 0,
\]
and $M'$ is invertible (triangular with unit diagonal). Undoing the
conjugation, $M = PM'P^{-1}$ is invertible and
$\operatorname{tr}(AM) = \operatorname{tr}(BM') = 0$: $M \in H_A$.
\end{solution}

\begin{solution}{exo:b2:linalg:7}
$\det(I + tA)$ is, by the permutation formula, a polynomial in $t$;
its constant term is $1$ ($t = 0$). Its $t$-coefficient: expand
$\det$ as an alternating form of the columns $e_j + t\,c_j(A)$; by
multilinearity, the terms linear in $t$ replace exactly one $e_j$ by
$c_j(A)$:
\[
\sum_{j} \det(e_1, \dots, c_j(A), \dots, e_n)
= \sum_j a_{jj} = \operatorname{tr} A ,
\]
(the determinant with all canonical columns except $c_j(A)$ in slot
$j$ picks the $j$-th diagonal entry). Hence the derivative at $0$ is
$\operatorname{tr} A$.

Let $g(t) = \det(\eu^{tA})$. The group property gives $g(s + t) =
g(s)g(t)$ (multiplicativity of $\det$), $g$ is differentiable, and
$g'(0) = \operatorname{tr} A$ by the above ($\eu^{tA} = I + tA +
O(t^2)$). A differentiable morphism $(\R, +) \to (\R^*, \times)$
satisfies $g' = g'(0)\,g$ (differentiate $g(s+t)$ in $s$ at $0$), so
$g(t) = \eu^{t\operatorname{tr} A}$ by the uniqueness of solutions
of $y' = cy$ with $y(0) = 1$ (Year 1 volume).
\end{solution}

\begin{solution}{exo:b2:linalg:8}
Let $v_k = (1, j^k, j^{2k})^{\mathsf T}$ for $k = 0, 1, 2$. Using $1
+ j + j^2 = 0$ and $j^3 = 1$:
\[
C v_k =
\begin{pmatrix}
a + b j^k + c j^{2k}\\
c + a j^k + b j^{2k}\\
b + c j^k + a j^{2k}
\end{pmatrix}
= (a + b j^k + c j^{2k})
\begin{pmatrix} 1\\ j^k\\ j^{2k}\end{pmatrix},
\]
(check the second row: $j^k(a + bj^k + cj^{2k}) = aj^k + bj^{2k} +
cj^{3k} = c + aj^k + bj^{2k}$). So $v_k$ is an eigenvector with
eigenvalue $\lambda_k = a + bj^k + cj^{2k}$. The $v_k$ form a basis
(Vandermonde of the distinct $1, j, j^2$), so $C$ is diagonalizable
with these eigenvalues and
\[
\det C = \lambda_0\lambda_1\lambda_2
= (a+b+c)(a + bj + cj^2)(a + bj^2 + cj).
\]
\end{solution}

\begin{solution}{exo:b2:linalg:9}
The map $\Theta \colon A \mapsto \operatorname{tr}(A\,\cdot)$ is
linear from $\mathcal{M}_n(K)$ to its dual, between spaces of equal
dimension $n^2$: injectivity suffices. If $\operatorname{tr}(AM) =
0$ for all $M$, take $M = E_{ji}$: $\operatorname{tr}(A E_{ji}) =
a_{ij} = 0$ for all $i, j$: $A = 0$. So $\Theta$ is an isomorphism.

Uniqueness of the trace (\cref{prop:b2:linalg:trace}): a form $t$
killing all commutators is $\operatorname{tr}(A\,\cdot)$ for some
$A$ with $\operatorname{tr}(A(MN - NM)) = 0$ for all $M, N$, i.e.\
$\operatorname{tr}((AM - MA)N) = 0$ for all $N$ (cyclicity), i.e.\
$AM = MA$ for all $M$ (injectivity of $\Theta$): $A$ commutes with
everything, hence is scalar ($A$ commutes with all $E_{ij}$ forces
off-diagonal entries $0$ and equal diagonal entries), so $t = c
\operatorname{tr}$.
\end{solution}

\begin{solution}{exo:b2:linalg:10}
Work over $\C$ (a real matrix is nilpotent iff it is as a complex
matrix: nilpotence is $u^n = 0$).

\emph{Step 1: $\operatorname{tr}(u^k) = 0$ for $k \geq 1$.} By
induction, $u^k v - v u^k = k\, u^k$: for $k = 1$ it is the
hypothesis; for the step,
\[
u^{k+1}v - vu^{k+1} = u^k(uv - vu) + (u^k v - v u^k)u
= u^{k+1} + k\,u^{k+1} .
\]
Taking traces: $0 = \operatorname{tr}(u^k v) -
\operatorname{tr}(vu^k) = k \operatorname{tr}(u^k)$, so
$\operatorname{tr}(u^k) = 0$.

\emph{Step 2: zero power traces imply nilpotence (over $\C$).} Let
$\lambda_1, \dots, \lambda_r$ be the distinct nonzero eigenvalues of
$u$ with multiplicities $m_1, \dots, m_r$ (in the characteristic
polynomial, which splits over $\C$ ---
\cref{ch:b2:reduction}). Power traces are $\operatorname{tr}(u^k) =
\sum_i m_i \lambda_i^k$ (trigonalize: the diagonal of a triangular
matrix's $k$-th power is the $k$-th powers). The system $\sum_i m_i
\lambda_i^k = 0$ for $k = 1, \dots, r$ is Vandermonde-invertible in
the unknowns $m_i\lambda_i$ (matrix $(\lambda_i^{k-1})$ times
diagonal $\lambda_i$, all $\lambda_i \neq 0$ distinct): every $m_i
\lambda_i = 0$, impossible with $m_i \geq 1$ unless $r = 0$. So $u$
has no nonzero eigenvalue: its characteristic polynomial is
$(-X)^n$, and Cayley--Hamilton (\cref{ch:b2:reduction}) gives $u^n
= 0$: nilpotent.
\end{solution}
